% Pursuing Stacks (cotes 134-2 to 134-8): archive cote 134-7, site page 47 (PDF page 48), typescript page 489.
% typescript. Draft compared: ps6.tex:2320-2387.
% AI-transcribed by Claude Opus 5.5 (claude-opus-5-5, adaptive thinking, effort high), Message Batches API, batch ps3 (msgbatch_01XdFmuVVBAcgc7EFxdtCqsj), 2026-11-12; not yet checked by eye.
\scanpage{134-7/47}
\mspage{\xout{(25)} 489}

if the \xout{xacx}\inserted[typed]{assumption} that (24) is a quasi-isomorphism should imply the same state-%
ment with $L(X_*)$ replaced by it's componentwise projective limit - name-%
ly \hand{that} (V) on p. 478 is, which we'll need of course in case $k=\underline{Z}$ for the
so-called ``linearization theorem''. \xout{Thexx} Thus we get three \xout{xelated}
isomorphism or quasi-isomorphism statement\hand{s,} concerning (23),(24) and
(V) in yesterday's notes, which are at any rate closely related, and
which \xout{Ixfeel} one hopes to be true, be\hand{c}ause this\inserted{seems} \deleted{is} needed for a schema-%
tization theory of homotopy types to work. But I s\hand{h}ould confess I have
not tried even to get any clue as to why this should be true, under the
only assumptions say that the components of $X_*$ should be smooth (and
possibly the Kan assumption ?), plus $X_0=e$ say.

Now to the second ingredient of the looked-for ``linearization theo-%
rem'', which previously was the first map in (w) p.454, or (V') on page
478, involving maps of \xout{either Z(Xx(k))}
\[ \text{either } \underline{Z}^{(X_{\hand{$*$}}(k))} \quad \text{or} \quad k^{(X_{\hand{$*$}}(k))} \]
into what was previously called $L(X_*)$, and which we \xout{now} would now
rather denote\deleted{s} by
\[ \underset{\longleftarrow}{\lim}\, L(X_*) \quad ( \simeq \Gamma_k\,\hat{}\,(M_{\hand{$*$}}) \ \text{if} \ X_{\hand{$*$}} = W(M_{\hand{$*$}}) ) \ . \]
We made sure that, unless $k=\underline{Z}$ , none of the two \over{\unsure{b}}{h}ad any chance to be
a quasi-isomorphism. The only positive thing that came out in this
direction was that the first one of these maps would induce an isomor-%
phism on $H_i$ in the critical dimension (namely $i=n+1$ in the n-connected
case). We now understand why, \xout{ex} for $k \neq \underline{Z}$ , we would not get any
actual quasiisomorphism - namely the ``correct'' naive linearization
which compares reasonably with $L(X_*)$ should not be relative to a
\uline{constant} ring such as $\underline{Z}$ or $k$ , but relative to $\underline{O}_k$ , via the map
($*$) (p. 487) giving rise now to
\margin{\fbox{$(X_*)$}}
\[ (25) \qquad \underline{O}_k^{(X_{\hand{$*$}})} \hand{$\longrightarrow$} L(X_*) \quad . \]
The more reasonable question now, making a good sense really for any
ground ring $k$, is whether this map (under the usual assumptions\inserted{say} on $X_*$)
is a quasi-isomorphism. I would'nt really bet it is, as I have some
doubts as to whether the \xout{co}homology sheaves of the first member (both
\hand{(}\inserted{members} of course being viewed as chain complexes of $\underline{O}_k$-Modules\inserted[typed]{or ``pro'' such}) are quasi-
coherent - but for the time being I am not sure either \hand{i}f those of
the second member are essentially co\hand{n}stant proobjects ! But even if
(25) is'nt a quasiisomorphism, it does behave like one for all practi-%
cal purposes of computing\inserted[typed]{quasi-coherent} cohomology\hand{,} it would seem, as \xout{xxxxxxxx}\inserted[typed]{this boils}
\inserted{down indeed to the isomorphisms (20) or (20').}

\margin{47}
